\section*{Bab \ref{ch:b2:coordination-complexes} --- Logam Transisi dan Kompleks Koordinasi}

\begin{solution}{exo:b2:coordination-complexes:1}
\ce{Ti^3+} $d^1$, \ce{V^3+} $d^2$, \ce{Cr^3+} $d^3$, \ce{Mn^2+} $d^5$, \ce{Fe^3+} $d^5$, \ce{Co^2+}
$d^7$, \ce{Ni^2+} $d^8$, \ce{Zn^2+} $d^{10}$.
\end{solution}

\begin{solution}{exo:b2:coordination-complexes:2}
Air 1, en 2, etanadioat 2, edta 6, tiosianat 1; tiosianat bersifat
ambidentat (melalui S atau N).
\end{solution}

\begin{solution}{exo:b2:coordination-complexes:3}
Tetraamintembaga(II) sulfat; kalium heksasianidoferat(III);
tetraamindikloridokobalt(III) klorida; tetrakarbonilnikel(0).
\end{solution}

\begin{solution}{exo:b2:coordination-complexes:4}
Linear; tetrahedral; bujur sangkar planar ($d^8$, deret ketiga); oktahedral.
\end{solution}

\begin{solution}{exo:b2:coordination-complexes:5}
Dua: cis (kedua klorida bertetangga) dan trans (berseberangan). Dalam
tetrahedron hanya akan ada satu.
\end{solution}

\begin{solution}{exo:b2:coordination-complexes:6}
\ce{Cr(CO)6} 18; \ce{Fe(CO)5} 18; \ce{Ni(CO)4} 18; \ce{V(CO)6}: $5 + 12 = 17$;
\ce{Mn2(CO)10} 18 per mangan. \ce{V(CO)6} melanggar aturan itu: senyawa ini
\omterm{def:b2:diatomic-mos:magnetism}{paramagnetik} dan mudah direduksi menjadi \ce{[V(CO)6]-}, yang mempunyai 18.
\end{solution}

\begin{solution}{exo:b2:coordination-complexes:7}
Ferosena: Fe(II) $d^6$ + 12 = 18. Kobaltosena: Co(II) $d^7$ + 12 = 19.
Kobaltosena mudah melepaskan elektron kelebihannya dan menghasilkan ion
kobaltosenium \ce{[Co(C5H5)2]+}, dengan 18.
\end{solution}

\begin{solution}{exo:b2:coordination-complexes:8}
Satu edta menggantikan enam molekul air: \ce{[Ni(H2O)6]^2+ + edta^4- -> [Ni(edta)]^2- +
6H2O}, tujuh partikel dari dua. \omterm{def:b2:reaction-free-energy:reaction-entropy}{Entropi reaksinya} sangat positif, dan $K$
sangat besar. Enam ligan terpisah menggantikan enam molekul air tanpa
penambahan jumlah partikel.
\end{solution}

\begin{solution}{exo:b2:coordination-complexes:9}
Nitrito-$\kappa N$ (\ce{Co-NO2}, isomer nitro, kuning) dan nitrito-$\kappa O$
(\ce{Co-O-N=O}, isomer nitrito, merah).
\end{solution}

\begin{solution}{exo:b2:coordination-complexes:10}
Tiga: isomer trans (akiral: ada bidang-bidang cermin) dan kedua enansiomer
isomer cis (kedua kelat dan kedua klorida melilit ke satu arah atau ke arah
yang lain).
\end{solution}

\begin{solution}{exo:b2:coordination-complexes:11}
Heksagon planar: B pada posisi 1,2 (orto), 1,3 (meta), 1,4 (para): tiga
isomer. Prisma trigonal: B sepanjang sisi segitiga, sepanjang rusuk tegak,
atau melintang pada muka persegi panjang: tiga. Oktahedron memberikan dua.
Tidak ditemukannya lebih dari dua isomer, meskipun banyak percobaan dan
beberapa senyawa telah dicoba, menyingkirkan kedua geometri yang lain ---
asalkan tidak ada isomer ketiga yang terlewat.
\end{solution}

\begin{solution}{exo:b2:coordination-complexes:12}
\ce{[RhCl(PPh3)3]}: Rh(I) $d^8$ + 4 × 2 = 16. \ce{[IrH(CO)(PPh3)3]}: Ir(I) $d^8$ + 5 × 2 =
18. Kompleks rodium, dengan 16 elektron, dapat mengikat dua lagi (adisi
oksidatif, \cref{ch:b2:catalytic-cycles}).
\end{solution}

\begin{solution}{pb:b2:coordination-complexes:1}
\textbf{1.} Klorida-klorida di luar \omterm{def:b2:coordination-complexes:coordination-sphere}{bola koordinasi}, yang bebas dalam larutan.
\textbf{2.} A: 0; B: 1; C dan D: 2.
\textbf{3.} A: $[\ce{Co(NH3)6}]^{3+}$ + 3 \ce{Cl-}: 4 ion; B: 3; C, D: 2. Hasilnya cocok.
\textbf{4.} +3.
\textbf{5.} $d^6$.
\textbf{6.} Enam dalam setiap senyawa: 6 \ce{NH3}; 5 \ce{NH3} + 1 Cl; 4 \ce{NH3} + 2 Cl.
\textbf{7.} \ce{[Co(NH3)6]Cl3}; \ce{[CoCl(NH3)5]Cl2}; \ce{[CoCl2(NH3)4]Cl} (dua kali).
\textbf{8.} Heksaaminkobalt(III) klorida; pentaaminkloridokobalt(III) klorida.
\textbf{9.} Tetraamindikloridokobalt(III).
\textbf{10.} Isomer geometri (cis--trans).
\textbf{11.} \ce{[CoCl3(NH3)3]}, sebuah kompleks netral: tidak ada ion, tidak
ada klorida yang langsung mengendap.
\textbf{12.} Dua: fac dan mer.
\textbf{13.} Cis: kedua Cl pada $90^\circ$; trans: pada $180^\circ$.
\textbf{14.} Tiga (posisi orto, meta, para pada heksagon).
\textbf{15.} Tiga.
\textbf{16.} Hal itu mendukung oktahedron, satu-satunya dari ketiga geometri
yang meramalkan dua.
\textbf{17.} Tidak: ketiadaan isomer ketiga adalah bukti negatif. Isomer
ketiga akan menyingkirkan oktahedron.
\textbf{18.} Isomer yang membentuk kelat: etanadioat yang bidentat hanya
dapat merentang dua posisi cis.
\textbf{19.} Tiga kelat en di sekeliling kobalt, masing-masing merentang dua
posisi cis.
\textbf{20.} Tidak keduanya: tidak ada bidang cermin dan tidak ada pusat inversi.
\textbf{21.} Dua enansiomer, $\Delta$ dan $\Lambda$.
\textbf{22.} Bahwa keenam ligan tersusun dalam tiga dimensi, seperti dalam
oktahedron; susunan planar akan menghasilkan ion yang akiral.
\textbf{23.} Cis: kiral (dua enansiomer). Trans: akiral.
\textbf{24.} Pemisahan itu menunjukkan bahwa aktivitas optik merupakan sifat
geometri molekul secara umum, bukan milik karbon, dan meneguhkan oktahedron.
\textbf{25.} Oktahedron meramalkan $\boldsymbol{2}$ isomer \ce{[CoCl2(NH3)4]+};
heksagon planar dan prisma trigonal meramalkan $\boldsymbol{3}$.
\end{solution}
